\documentclass[a4paper,11pt]{article}

\usepackage{amsmath}
\usepackage{url}
\usepackage{hyperref}
\usepackage{enumitem}

\title{Milestone 1: Technical Understanding and Research Framing}
\author{Christian Percival}
\date{}

\begin{document}

\maketitle

\section*{Topic}

\begin{itemize}
    \item Seminar topic: Handling execution overruns in hard real-time control systems.
    \item Main reference paper:
    \begin{itemize}
        \item M. Caccamo, G. Buttazzo, and L. Sha, ``Handling Execution Overruns in Hard Real-Time Control Systems,'' IEEE Transactions on Computers, 2002.
    \end{itemize}
    \item Official background reference:
    \begin{itemize}
        \item G. Buttazzo, \textit{Hard Real-Time Computing Systems: Predictable Scheduling Algorithms and Applications}, 4th ed., Springer, 2024.
    \end{itemize}
\end{itemize}

\section*{Problem Statement}

\begin{itemize}
    \item Hard real-time systems must complete tasks before their deadlines.
    \item A result that is correct but produced too late can still be wrong in a hard real-time system.
    \item Classical scheduling often uses the worst-case execution time (WCET).
    \item WCET-based scheduling is safe, but often inefficient.
    \item Many tasks usually execute much faster than their WCET.
    \item The problem is how to use the processor efficiently in normal cases while still remaining safe when worst-case execution happens.
\end{itemize}

\section*{What is an Execution Overrun?}

\begin{itemize}
    \item An execution overrun happens when a task executes longer than its expected or reserved execution time.
    \item It does not mean that the task finishes earlier than WCET.
    \item Finishing earlier than WCET only means that reserved processor time may be unused.
    \item In the main paper, the normal computation time is denoted as:
    \[
        c_i^n
    \]
    \item The worst-case execution time is denoted as:
    \[
        WCET_i
    \]
    \item An overrun occurs when:
    \[
        c_i > c_i^n
    \]
    \item The overrun amount is:
    \[
        c_i - c_i^n
    \]
\end{itemize}

\section*{Motivation from Visual Tracking}

\begin{itemize}
    \item The main paper motivates the problem using visual tracking.
    \item A vision task may normally search only a small predicted image window.
    \item If the object is found in this small window, the task finishes quickly.
    \item If the object is not found, the search area must be expanded.
    \item Searching a larger area takes more computation time.
    \item This creates a task with:
    \begin{itemize}
        \item short normal execution time,
        \item occasional long execution time,
        \item large WCET.
    \end{itemize}
    \item Similar situations occur in:
    \begin{itemize}
        \item robotics,
        \item radar tracking,
        \item autonomous systems,
        \item sensor-based control systems.
    \end{itemize}
\end{itemize}

\section*{Core Technical Idea}

\begin{itemize}
    \item Do not always schedule tasks only according to WCET.
    \item Use normal execution time to allow higher control frequencies.
    \item Use WCET information to guarantee safety.
    \item Use bandwidth reservation to prevent one task from using too much CPU time.
    \item Use deadline postponement to handle extra execution time.
    \item Use server-based scheduling to isolate the effect of overruns.
\end{itemize}

\section*{Important Terms}

\begin{itemize}
    \item \textbf{Hard real-time task:} a task where missing a deadline is unacceptable.
    \item \textbf{WCET:} worst-case execution time.
    \item \textbf{Normal computation time:} typical runtime of the task.
    \item \textbf{Execution overrun:} task exceeds expected or reserved execution time.
    \item \textbf{EDF:} Earliest Deadline First scheduling.
    \item \textbf{CBS:} Constant Bandwidth Server.
    \item \textbf{CBShd:} CBS extension for hard-deadline overrun handling.
    \item \textbf{PLI:} Performance Loss Index.
    \item \textbf{Temporal isolation:} preventing one task from disturbing other tasks too much.
\end{itemize}

\section*{Mathematical Model}

\begin{itemize}
    \item Each control task can be modeled as:
    \[
        \tau_i(WCET_i, c_i^n, f_i^{min}, \Delta J_i(f_i))
    \]
    \item $WCET_i$: worst-case execution time.
    \item $c_i^n$: normal computation time.
    \item $f_i^{min}$: minimum allowed frequency.
    \item $\Delta J_i(f_i)$: performance loss index.
\end{itemize}

\section*{Utilization}

\begin{itemize}
    \item Task utilization can be written as:
    \[
        U_i = C_i f_i
    \]
    \item Total utilization:
    \[
        U = \sum_{i=1}^{n} U_i
    \]
    \item For EDF on one processor, the task set is usually schedulable if:
    \[
        \sum_{i=1}^{n} U_i \leq 1
    \]
\end{itemize}

\section*{Performance Loss Index}

\begin{itemize}
    \item The paper models performance loss as:
    \[
        \Delta J_i(f_i) = \alpha_i e^{-\beta_i f_i}
    \]
    \item $f_i$: task frequency.
    \item $\alpha_i$: magnitude coefficient.
    \item $\beta_i$: decay rate.
    \item Higher frequency usually reduces performance loss.
    \item Higher frequency also increases CPU usage.
    \item The total performance loss is:
    \[
        \Delta J(f_1,\ldots,f_n) = \sum_i w_i \Delta J_i(f_i)
    \]
\end{itemize}

\section*{CBS and CBShd}

\begin{itemize}
    \item CBS assigns each task a server budget and server period:
    \[
        (Q_s, T_s)
    \]
    \item The server bandwidth is:
    \[
        U_s = \frac{Q_s}{T_s}
    \]
    \item CBS limits how much processor time a task can consume.
    \item If the task exhausts its budget, the deadline is postponed.
    \item CBShd improves CBS for hard-deadline tasks.
    \item CBShd avoids postponing deadlines too far when the remaining overrun is small.
\end{itemize}

\section*{Assumptions}

\begin{itemize}
    \item Tasks are periodic control tasks.
    \item WCET is known or can be estimated.
    \item Normal execution time is known or measured.
    \item Minimum safe task frequency is known.
    \item Overruns are bounded by WCET.
    \item The system uses EDF or a compatible server-based scheduler.
    \item Remaining computation time can be estimated.
\end{itemize}

\section*{Limitations}

\begin{itemize}
    \item WCET is difficult to estimate on modern processors.
    \item Normal execution time depends on input data.
    \item Server management introduces runtime overhead.
    \item Smaller budgets improve precision but increase scheduling overhead.
    \item Implementation is easier in a real-time kernel than in a normal desktop operating system.
\end{itemize}

\section*{Open Questions}

\begin{itemize}
    \item How accurately can WCET be measured in a student implementation?
    \item Should the final paper focus more on CBS/CBShd theory or the implementation example?
    \item Is a Python simulation acceptable as an implementation demonstration?
    \item How should the implementation distinguish normal cases, overruns, and worst cases?
    \item How much mathematical detail should be included in the final 5-page paper?
\end{itemize}

\end{document}